\section{OS Forensics}

Operating System Forensics covers the extraction and analysis of architectural artifacts left by the OS kernel, system services, and user-space management applications.

\subsection{Volatile Memory Forensics}

Memory forensics involves the capture and programmatic triage of volatile RAM. 

\subsubsection*{Capture Lifecycle}
Since RAM is always changing due to how the operating system manages it and runs programs, investigators need to collect data from it using special methods that are fast and don't use much space.

\begin{itemize}
    \item \textbf{Live Acquisition Frameworks:} Utilities running within user or kernel space (e.g., \texttt{LiME} for Linux, \texttt{FTK Imager Lite} or \texttt{DumpIt} for Windows) copy the physical memory space into a raw dump format.
    \item \textbf{Hardware-Level Interception:} Using Direct Memory Access (DMA) attacks over hardware interfaces (e.g., Thunderbolt) to completely bypass OS-level instrumentation and rootkits.
\end{itemize}

\noindent
Once a volatile memory dump is secured, investigators analyze the binary structure using frameworks like \textbf{Volatility} to extract runtime artifacts.
\medskip

\noindent 
Key volatile memory artifacts commonly extracted during analysis include:
\begin{itemize}
    \item \textbf{Active Process Trees}: maps active and terminated process structures
    \item \textbf{Network Sockets}: shows open connections and listening ports, useful to link to C2 infrastructure.
    \item \textbf{Encryption Keys}: recovers in-memory keys and session tokens.
    \item \textbf{Cached File Handles}: lists files opened by processes.
\end{itemize}

\subsection{Persistent OS Artifacts and Configuration Repositories}

Persistent artifacts remain intact across reboots within specific file system locations. 

\subsubsection*{Windows Registry}
The registry acts as a centralized database for OS configurations, user preferences, and hardware properties. It is structured hierarchically into binary files called \textbf{hives}.
\begin{itemize}
    \item \textbf{SYSTEM / SAM Hives:} Contains machine security configurations, hardware definitions, and localized user password hashes.
    \item \textbf{NTUSER.DAT / USRCLASS.DAT:} User-specific hives located within individual profile directories. They hold execution histories, recently accessed file lists (MRU), and desktop environment settings.
\end{itemize}

\subsubsection*{User Activity and Execution Evidence}
To establish user intent and reconstruct an event timeline, investigators query specialized system tracking structures:
\begin{itemize}
    \item \textbf{Prefetch Files (\texttt{.pf}):} Designed to optimize application launch times.\\
    Prefetch files capture the application execution history, execution timestamps, volume serial numbers, and a list of dependent DLL files.
    \item \textbf{Shimcache (AppNameCache):} A kernel level caching mechanism used to track application compatibility properties.\\
    It records file names, paths, sizes, and the last modification time of executed binaries.
    \item \textbf{Amcache:} A hive file that records detailed metadata regarding application installations and executions, capturing SHA-1 file hashes and path strings even if the underlying binary has been deleted from disk.
\end{itemize}

\subsubsection*{System Event Logs}
Centralized logging daemons record application, security, and administrative operations. In Windows environments, these are managed as binary \textbf{EVTX} files.
\begin{itemize}
    \item \textbf{Security Log (Event ID 4624/4625):} Logs successful and failed authentication attempts, recording the logon type (e.g., interactive, network, RDP), target username, and source IP address.
    \item \textbf{System/Application Logs:} Captures service state modifications, driver errors, hardware connections, and critical software crashes.
\end{itemize}

\subsection{Linux and Unix Artifacts}

Unlike the centralized registry structure of Windows, Linux and Unix-like put everything in a file.\\


\subsubsection*{Pluggable Authentication Modules (PAM)}
Authentication mechanisms and account configurations are managed globally within the \texttt{/etc} configuration directory.
\begin{itemize}
    \item \textbf{\texttt{/etc/passwd}:} Contains essential user account information, listing usernames, user IDs (UID), group IDs (GID), home directory paths, and the default command-line shell interface.
    \item \textbf{\texttt{/etc/shadow}:} A restricted system file that stores encrypted user password hashes alongside account aging parameters. Access requires root privileges, making its integrity a primary focus during post-compromise triage.
    \item \textbf{\texttt{/etc/pam.d/}:} Houses configuration rules for the PAM framework. Investigators audit these files to identify unauthorized modifications, such as backdoor modules allowing password bypass or unauthorized privilege elevation.
\end{itemize}

\subsubsection*{Chronological Command Execution and Shell History}
Reconstructing user interaction timelines requires parsing volatile and semi-persistent command-line history logs.
\begin{itemize}
    \item \textbf{\texttt{.bash\_history} / \texttt{.zsh\_history}:} Hidden files located within each user's home directory that record executed commands sequentially.
\end{itemize}

\noindent
By default, command history files are only updated upon clean shell termination and do not inherently include execution timestamps. Furthermore, attackers can easily disable or manipulate history tracking via environment variables (e.g., \texttt{unset HISTFILE} or setting \texttt{HISTCONTROL=ignorespace}).

\subsubsection{Persistence Mechanisms}
Malware achieves persistence on Linux systems by integrating malicious execution paths into automation hooks and system initialization frameworks.

\begin{table}[H]
\centering
\small
\begin{tabular}{|p{3.5cm}|p{8.5cm}|}
\hline
\textbf{Persistence Mechanism} & \textbf{Forensic Inspection Paths and Behavior} \\
\hline
\textbf{Cron Jobs} & Auditing \texttt{/etc/crontab}, \texttt{/etc/cron.*/}, and individual user crontabs via \texttt{/var/spool/cron/}. These files schedule periodic background scripts and command executions. \\
\hline
\textbf{Systemd Daemons} & Inspecting unit files within \texttt{/etc/systemd/system/} and \texttt{/lib/systemd/system/}. Systemd manages background services, allowing malicious binaries to start automatically at boot with high privileges. \\
\hline
\textbf{Shell Profiles} & Reviewing global and local rc files (\texttt{/etc/profile}, \texttt{\textasciitilde/.bashrc}, \texttt{\textasciitilde/.zshrc}). Malicious commands embedded here execute automatically every time a targeted user spawns a new shell session. \\
\hline
\end{tabular}
\caption{Primary Linux persistence mechanism vectors and forensic audit paths.}
\end{table}

\subsubsection*{Centralized Logging}
The standard logging infrastructure on Unix-like environments aggregates system, kernel, and application messages into centralized, plain-text files located within the \texttt{/var/log/} directory.

\begin{itemize}
    \item \textbf{\texttt{/var/log/auth.log} (or \texttt{secure}):} Records authentication states, capturing remote connection traces (e.g., SSH logons), privilege elevations via \texttt{sudo}, and local account modifications.
    \item \textbf{\texttt{/var/log/syslog} (or \texttt{messages}):} Serves as a global repository for non-critical system components, including daemon start/stop cycles, device attachments, and runtime application errors.
    \item \textbf{Log Rotation and Ingestion:} System logs are periodically rotated (compressed into \texttt{.gz} format) via the \texttt{logrotate} facility. This behavior leaves historical compressed log fragments behind, which can yield vital historical evidence during deep disk analysis.
\end{itemize}

\subsubsection*{Binary State Accounting}
To optimize access speeds and prevent system logs from inflating rapidly, connection states are recorded within structured, fixed-size binary data files that cannot be edited via standard text editors.
\begin{itemize}
    \item \textbf{\texttt{utmp}:} Maintains real-time tracking of currently authenticated users, their bound terminal endpoints, and their source IP addresses.
    \item \textbf{\texttt{wtmp}:} Acts as a historical ledger, recording an append-only archive of all successful user logins, logouts, and system boot/reboot timestamps.
    \item \textbf{\texttt{btmp}:} Records historical data for failed authentication attempts, serving as a primary telemetry source for identifying automated brute-force attacks or scanning patterns.
\end{itemize}

\noindent
These binary files are queried using specialized operating system utilities such as \texttt{last}, \texttt{lastb}, or \texttt{who}, bypassable only if an attacker manually zeros out or deletes the underlying data structures.

\subsection{Volatile Analysis Frameworks}

When data is cleared or unallocated during nominal deletion on a storage drive, the information is not instantly erased from central memory.\\
A data segment or application payload remains intact in RAM. 

\subsubsection*{Volatility Architecture}
The \textbf{Volatility} framework is an open-source, multi-platform forensic suite

\begin{itemize}
    \item \textbf{Volatility 2:} Legacy versions require investigators to specify a rigid OS target profile.\\ This profile makes difficult to analyze memory dumps from unknown or custom builds.
    \item \textbf{Volatility 3:} Modern versions migrate away from rigid system profiles, utilizing automated, structured JSON symbol tables that map kernel offsets dynamically for specific kernel builds.
\end{itemize}

\noindent
Volatility modules help investigators quickly extract runtime traces and detect suspicious behavior:
\begin{itemize}
    \item \textbf{Behavioral analysis:} identify suspicious processes and network connections.
    \item \textbf{Rootkit detection:} search for kernel manipulation and hidden drivers.
    \item \textbf{Data extraction:} recover ephemeral keys, credentials, and in-memory application data.
\end{itemize}

